\chapter{Equivalent Probability Measures and Change of Measure}

In finance (Girsanov Theorem), we frequently move between the ``historical'' or ``physical'' measure $\mathbb{P}$ and a ``risk-neutral'' measure $\mathbb{Q}$.

\section{Radon-Nikodym Derivative}

\begin{definition}[Radon-Nikodym Density]
Let $\mathbb{P}$ be a reference probability measure. Let $Z$ be a non-negative random variable ($Z \geq 0$ a.s.) with $\mathbb{E}_{\mathbb{P}}[Z] = 1$.
We can define a new probability measure $\mathbb{Q}$ by weighing the outcomes by $Z$:
\[
\mathbb{Q}(A) := \int_A Z(\omega) \, d\mathbb{P}(\omega) = \mathbb{E}_{\mathbb{P}}[ \mathbb{1}_A Z ], \quad \forall A \in \mathcal{F}.
\]
The variable $Z$ is called the \textbf{Radon-Nikodym derivative} (or density) of $\mathbb{Q}$ with respect to $\mathbb{P}$. We write:
\[
Z = \frac{d\mathbb{Q}}{d\mathbb{P}}.
\]
\end{definition}

\begin{remark}[Notation: Indicator Function]
The symbol $\mathbb{1}_A$ is the \textbf{Indicator Function}.
\[
\mathbb{1}_A(\omega) = \begin{cases} 
1 & \text{if } \omega \in A \\
0 & \text{if } \omega \notin A 
\end{cases}
\]
So $\mathbb{E}_{\mathbb{P}}[\mathbb{1}_A Z]$ is just a fancy way of writing "The integral of $Z$ only over the set $A$".
\end{remark}

\paragraph{Intuition: Multiple Probability Measures?}
How can we have two probabilities $\mathbb{P}$ and $\mathbb{Q}$?
\begin{itemize}
    \item Think of a \textbf{Coin Toss}.
    \item $\Omega = \{H, T\}$. The outcomes are fixed.
    \item Measure $\mathbb{P}$ (Fair Coin): $\mathbb{P}(H) = 0.5, \mathbb{P}(T) = 0.5$.
    \item Measure $\mathbb{Q}$ (Biased Coin): $\mathbb{Q}(H) = 0.6, \mathbb{Q}(T) = 0.4$.
    \item We just "re-weighted" the probabilities. The universe $\Omega$ didn't change, only our assignment of likelihoods.
    \item The factor $Z$ tells us how to re-weight: $Z(H) = \frac{0.6}{0.5} = 1.2$, $Z(T) = \frac{0.4}{0.5} = 0.8$.
\end{itemize}

\subsection{Relationship between Expectations}
For any random variable $X$:
- $X \in L^1(\mathbb{Q})$ iff $XZ \in L^1(\mathbb{P})$.
- The expectation transforms as:
\[
\mathbb{E}_{\mathbb{Q}}[X] = \mathbb{E}_{\mathbb{P}}[X Z].
\]

\paragraph{Intuition: Why multiply by \texorpdfstring{$Z$}{Z}?}
When we compute an average under $\mathbb{Q}$, we should weight outcomes by their $\mathbb{Q}$-probability. However, if we want to use the $\mathbb{P}$-probability (because we are computing $\mathbb{E}_{\mathbb{P}}$), we must adjust the weight of each outcome $\omega$.
\[
\text{Weight in } \mathbb{Q} = Z(\omega) \times \text{Weight in } \mathbb{P}.
\]
So, to compute the $\mathbb{Q}$-average of $X$, we take the $\mathbb{P}$-average of the \textbf{re-weighted variable} $X \cdot Z$.
\begin{itemize}
    \item If $Z(\omega) > 1$, outcome $\omega$ is \textbf{more likely} under $\mathbb{Q}$ than $\mathbb{P}$, so we give $X(\omega)$ \textbf{more weight}.
    \item If $Z(\omega) < 1$, outcome $\omega$ is \textbf{less likely} under $\mathbb{Q}$, so we give $X(\omega)$ \textbf{less weight}.
\end{itemize}

\subsection{Example: Gaussian Shift (Essential for Finance)}
This is the building block of the Girsanov Theorem.
\begin{itemize}
    \item Let $X \sim \mathcal{N}(0, 1)$ under $\mathbb{P}$. So $f_{\mathbb{P}}(x) = \frac{1}{\sqrt{2\pi}} e^{-x^2/2}$.
    \item We want to move to a measure $\mathbb{Q}$ where $X \sim \mathcal{N}(\mu, 1)$.
    \item The target density is $f_{\mathbb{Q}}(x) = \frac{1}{\sqrt{2\pi}} e^{-(x-\mu)^2/2}$.
    \item The Radon-Nikodym derivative is the ratio of densities:
    \[
    Z(\omega) = \frac{f_{\mathbb{Q}}(X(\omega))}{f_{\mathbb{P}}(X(\omega))} = \frac{e^{-(X-\mu)^2/2}}{e^{-X^2/2}} = e^{\mu X - \frac{1}{2}\mu^2}.
    \]
    \item If we use this $Z$ to re-weight probabilities, $X$ behaves like it has mean $\mu$ instead of $0$.
\end{itemize}

\section{Absolute Continuity and Equivalence}

\begin{definition}[Absolute Continuity]
A measure $\mathbb{Q}$ is \textbf{absolutely continuous} with respect to $\mathbb{P}$ (denoted $\mathbb{Q} \ll \mathbb{P}$) if $\mathbb{P}(A) = 0 \implies \mathbb{Q}(A) = 0$.
Intuitively, ``impossible'' events under $\mathbb{P}$ must remain ``impossible'' under $\mathbb{Q}$.
\end{definition}

\begin{definition}[Equivalent Measures]
Measures $\mathbb{P}$ and $\mathbb{Q}$ are \textbf{equivalent} ($\mathbb{P} \sim \mathbb{Q}$) if $\mathbb{Q} \ll \mathbb{P}$ and $\mathbb{P} \ll \mathbb{Q}$. They agree on which events are possible (have probability $> 0$) and which are impossible.
\end{definition}

\begin{theorem}[Radon-Nikodym Theorem]
If $\mathbb{Q} \ll \mathbb{P}$, there exists a unique measurable function $Z \geq 0$ such that $d\mathbb{Q} = Z d\mathbb{P}$.
If $\mathbb{Q} \sim \mathbb{P}$, then strictly $Z > 0$ almost surely.
\end{theorem}

\section{Conditional Bayes Formula/Kallianpur-Striebel}
How does conditional expectation change when we change the measure? This is crucial for pricing (computing $\mathbb{E}^{\mathbb{Q}}[ \dots | \mathcal{F}_t ]$).

\begin{theorem}[Conditional Bayes Formula]
Let $\mathbb{Q} = Z \cdot \mathbb{P}$. Let $\mathcal{G} \subset \mathcal{F}$ be a sub-$\sigma$-algebra. For any $X \in L^1(\mathbb{Q})$:
\[
\mathbb{E}_{\mathbb{Q}}[X | \mathcal{G}] = \frac{\mathbb{E}_{\mathbb{P}}[X Z | \mathcal{G}]}{\mathbb{E}_{\mathbb{P}}[Z | \mathcal{G}]}.
\]
\end{theorem}
\textbf{Intuition}: The ``local'' density on the subspace given by $\mathcal{G}$ is represented by the conditional expectation of $Z$. We must re-normalize by this local density.
